\subsection{仿射韦伊群}

对于 \(\alpha \in \mathbb{R}\) 和 \(k\in \mathbf{Z}\)，令 \(\mathrm{L}_{\alpha ,k}\) 为 E 中由下式定义的超平面：

\[
\mathrm{L} _ {\alpha , k} = \{x \in \mathrm{E} | \langle \alpha , x \rangle = k \}
\]

并令 \(s_{\alpha, k}\) 为关于 \(L_{\alpha, k}\) 的正交反射。于是

\[
s _ {\alpha , k} (x) = x - (\langle \alpha , x \rangle - k) \alpha^ {\prime} = s _ {\alpha , 0} (x) + k \alpha^ {\prime}
\]

对所有 \(x \in E\) 成立。换言之，

\[
s _ {\alpha , k} = t (k \alpha^ {\prime}) \circ s _ {\alpha}\tag{1}
\]

其中 \(s_{\alpha}\) 是关于超平面 \(L_{\alpha} = L_{\alpha,0}\) 的正交反射，即与根 \(\alpha\) 相关联的反射。

公式 (1) 表明，\(s_{\alpha, k}\) 不依赖于所选的标量积。

定义 1. 称由反射 \(s_{\alpha,k}\)（其中 \(\alpha \in R\) 且 \(k \in Z\)）生成的 E 的仿射变换群为根系 R 的仿射韦伊群，记为 \(\mathrm{W}_{a}(\mathrm{R})\)（或简记为 \(W_{a}\)）。

命题 1. 群 \(\mathrm{W}_a\) 是 \(\mathrm{W}\) 与 \(\mathrm{Q}\) 的半直积。

由于 W 由反射 \(s_{\alpha}\) 生成，所以 W 包含于 \(W_{a}\) 中。另一方面，\(t(\alpha^{\prime}) = s_{\alpha,1} \circ s_{\alpha}\)（当 \(\alpha \in R\) 时），这说明 \(Q \subseteq W_{a}\)。

由于 W 保持 \(Q(R^{\prime})\) 不变（\(\S\) 1，第 9 号），由 W 和 Q 生成的仿射变换群 G 是 W 与 Q 的半直积。由上可知 \(G \subseteq W_{a}\)，且根据 (1) 有 \(s_{\alpha,k} \in G\)（对所有 \(\alpha \in R\) 和 \(k \in Z\)）。因此 \(W_{a} = G\)。

命题 2. 群 \(W_{a}\) 赋予离散拓扑后在 E 上适当作用，并置换超平面 \(L_{\alpha,k}\)（其中 \(\alpha \in R\) 且 \(k \in Z\)）。

由于 \(Q(R^{\prime})\) 是 \(V^{*}\) 的离散子群，群 Q 在 E 上适当作用。因此 \(W_{a} = W.Q\) 也如此，因为 W 是有限群。此外，对于 \(\alpha, \beta \in R\) 和 \(k \in Z\)：

\[
s _ {\beta} (\mathrm{L} _ {\alpha , k}) = \mathrm{L} _ {\gamma , k} \text {with} \gamma = s _ {\beta} (\alpha) \in \mathrm{R},
\]

\[
t (\beta^ {\prime}) \left(\mathrm{L} _ {\alpha , k}\right) = \mathrm{L} _ {\alpha , k + n (\alpha , \beta)},
\]

其中 \(n(\alpha, \beta) = \langle \beta^{\prime}, \alpha \rangle\) 是整数，故得第二个断言。

因此可以将第五章第 3 节关于 \(W_{a}\) 在 E 上作用的结果应用于此。为避免与 W 在 \(V^{*}\) 中的室混淆，我们把由超平面系 \(L_{\alpha,k}\)（其中 \(\alpha \in R\) 且 \(k \in Z\)）在 E 中确定的室称为小室。于是 \(W_{a}\) 在小室集合上单纯传递地作用，而小室的闭包是 \(W_{a}\) 在 E 上作用的基本域（第五章第 3 节第 2 号定理 1 及第 3 号定理 2）。显然，W 群可与 \(U(W_{a})\)（\(W_{a}\) 在 \(V^{*}\) 的正交群中的典范像）相识别（参见第五章第 3 节第 6 号）。由此 \(W_{a}\) 是本质的（第五章第 3 节第 7 号），且 \(W_{a}\) 不可约当且仅当根系 R 不可约（第 1 节第 2 号命题 5 的推论）。若 R 不可约，每个小室都是开单纯形（第五章第 3 节第 9 号命题 8）。一般情况下，仿射空间 E 的典范乘积分解（第五章第 3 节第 8 号）对应于 R 分解为不可约分支。特别地，小室是开单纯形的乘积。

还要注意，第五章第 3 节第 2 号定理 1 的推论表明，\(s_{\alpha,k}\) 是 \(\mathbf{W}_a\) 中唯一的反射。

